\section{Discusión}

\subsection{Implicaciones para Misiones Reales}

Resulta revelador observar cómo un marco que combina generación y aprendizaje multi-agente deja de ser un ejercicio académico para convertirse en habilitador concreto de nuevas capacidades operacionales. 

La integración propuesta facilita el paso de operaciones supervisadas a misiones con supervisión mínima, especialmente en escenarios de servicado donde la latencia de comunicaciones resulta prohibitiva. Inspirados en validaciones de arquitecturas de referencia en plataformas reales \cite{Basciani2026_RA}, el framework permite a un servicer generar, evaluar y ejecutar planes completos de inspección o reparación con intervención humana solo en casos de alta criticidad. 

Aplicaciones prácticas como Li2025 AIAgents demuestran que agentes con bucles ReAct ya comienzan a operar de forma efectiva en entornos ground/space, sugiriendo que nuestra extensión híbrida podría acelerar la adopción en constelaciones comerciales y misiones de exploración profunda \cite{Li2025_AIAgents}.

\begin{table}[t]
\centering
\caption{Trade-offs arquitectónicos del marco propuesto}
\label{tab:tradeoffs}
\begin{tabular}{lccc}
\hline
Aspecto & Ventaja Propuesta & Desventaja & Comparación con Alternativas \\
\hline
Autonomía & Alta (híbrida) & Mayor complejidad & Superior a MAPE-K clásico \\
Eficiencia Computacional & Media-Alta & Consumo en inferencia LLM & Mejor que LLM puro \\
Verificabilidad & Alta (SCP + escudos) & Overhead de monitoreo & Superior a RL puro \\
Escalabilidad Multi-Agente & Alta & Entrenamiento costoso & Comparable a MARL federado \\
\hline
\end{tabular}
\end{table}

\subsection{Limitaciones Técnicas y Éticas}

Uno de los desafíos persistentes al madurar estos sistemas radica en distinguir entre lo técnicamente factible y lo operacionalmente responsable. 

Aunque el marco mejora significativamente métricas de desempeño, todavía exhibe vulnerabilidades frente a distribuciones de datos fuera del entrenamiento y posibles ataques adversariales sobre los componentes generativos. La dependencia de modelos de lenguaje grandes introduce además riesgos de opacidad en la trazabilidad de decisiones, algo que \cite{Li2025_AIAgents} ya señala como preocupación creciente en operaciones satelitales.

Desde el punto de vista ético, la delegación de autoridad decisoria a agentes autónomos en entornos orbitales congestionados plantea interrogantes sobre responsabilidad en caso de fallos y sobre la proliferación de capacidades duales. No obstante, cabe matizar que estas limitaciones no son exclusivas de nuestro enfoque, sino inherentes al estado actual de la autonomía espacial.

\subsection{Comparación con Estado del Arte}

Lejos de ser un asunto resuelto, la trusted autonomy en operaciones espaciales sigue evolucionando entre promesas y realidades. 

Nuestro marco se distingue por su integración explícita de razonamiento generativo con optimización verificable, superando en varios indicadores a enfoques puramente reactivos o puramente generativos. Comparado con revisiones comprehensivas del campo \cite{Thangavel2024}, ofrece una contribución concreta al cerrar la brecha entre planificación cognitiva y control seguro, algo que muchas arquitecturas DSS aún no resuelven satisfactoriamente.

En nuestra experiencia, esta hibridación no solo mejora métricas cuantitativas sino que aporta una flexibilidad operativa difícil de conseguir con métodos tradicionales. Claro está que trabajos como \cite{Basciani2026_RA} establecen un listón alto en adaptabilidad real, y nuestro enfoque busca precisamente alinearse con esa línea de madurez tecnológica.

En definitiva, más que reemplazar enfoques existentes, el marco propuesto busca complementarlos, proporcionando un camino pragmático hacia mayor autonomía sin sacrificar las garantías que las misiones espaciales exigen.